%\subsection{Cluster-Level Matrix Unit Integration}
%Virgo~\cite{kim2025virgo} explored a complementary direction by physically disaggregating the matrix unit from individual SIMT cores and placing a larger systolic array at the cluster level. By drawing operands directly from shared memory, Virgo increases operation granularity, improves data reuse, and reduces instruction-processing energy. Evaluations show substantial power and energy reductions compared with core-coupled Ampere- and Hopper-style designs. In contrast to Virgo’s approach of moving the matrix unit outside the SIMT core, FourShadow remains within the core-coupled, inner-product based Tensor Core model and explores its internal design space, rather than changing the unit’s placement in the hierarchy.

\section{Conclusion}

In this work, we presented \textit{FourShadow}, an open-source, synthesizable exploration framework built on the RISC-V Vortex GPGPU to systematically analyze SIMT Tensor Core microarchitectural scaling. By sweeping architectural parameters across ASIC synthesis (ASAP7 7\,nm), FPGA prototyping (AMD Alveo U55C), and cycle-accurate SIMT simulation, we demystified critical hardware trade-offs across the SIMT hierarchy.

%\begin{enumerate}
%    \item \textit{SIMT Hierarchy Scaling:} Tensor Core scaling is non-monotonic with warp width. A 16-thread configuration ($NT=16$) consistently yields the optimal area--performance Pareto frontier, achieving up to $27.46\times$ speedup over scalar execution for FP8. Furthermore, scaling issue width beyond $IW=4$ yields diminishing returns as non-TCU execution and operand delivery overheads dominate.
%    \item \textit{Datapath and Operand Staging:} Expanding the arithmetic datapath to a $2K$-wide FEDP reduces execution cycles by 28--30\%, absorbing minor pipeline latency increases. Conversely, offloading small control descriptors from registers into TCU-local SRAM introduces an unnecessary 1.2--6.5\% cycle penalty due to added synchronization and memory traffic.
%    \item \textit{Sparsity and Workload Mapping:} Achievable speedups from 2:4 structured sparsity ($2.21$--$2.35\times$ for WGMMA) depend on whether physical compressed mappings preserve full FEDP array utilization. When aligned with application memory intensity, combined WGMMA and sparse acceleration achieve end-to-end speedups of $9.6\times$ and $2.3\times$ for Llama~2 prefill and decode, respectively, and up to $26.82\times$ on compute-heavy NeRF workloads.
%\end{enumerate}

Beyond core-level insights, FourShadow serves as a robust substrate for holistic, open-source microarchitectural design space exploration (DSE). Moving forward, we plan to extend the framework in two key directions: (1)~integrating Distributed Shared Memory (DSMEM) mechanisms to enable multi-core warpgroups to exchange and broadcast matrix tiles directly across local SRAMs without hitting L2 cache bottlenecks and (2)~broadening microarchitectural DSE to explore emerging low-precision formats (e.g., NVFP4, MXFP), novel warp scheduling heuristics, and custom hardware--software co-design extensions.